\chapter{Composition des routes, des registres et des signatures}
\label{chap:composicion-especies-rev6}
\index{composition des routes}
\index{composition des registres}
\index{méson!expression compositionnelle}
\index{baryon!expression compositionnelle}

Une espèce composée ne peut entrer dans la loi des masses comme un nom auquel on attribue rétrospectivement une ligne de cinq entiers. Avant la signature doit exister un objet susceptible d'être composé. L'atlas du \cref{chap:atlas-genealogia-masas} fournit l'alphabet fini ; l'extension survivante et le registre du \cref{chap:extension-ruta-caracter} fournissent la mémoire susceptible de composition. Ce chapitre a pour objet de construire le niveau intermédiaire sans consulter de masses cibles :
\[
\begin{gathered}
 \text{configuration compositionnelle typée}
 \longrightarrow
 \text{registres complets de ses constituants}\\
 \longrightarrow
 \text{registre composé}
 \longrightarrow
 \ZZ^5.
\end{gathered}
\]

L’ordre des flèches fait partie du résultat. Les expressions se forment avant de recevoir une réalisation physique ; les enregistrements se composent avant l’extraction de la signature ; la signature n’est évaluée par le caractère qu’après la composition de la mémoire. Ainsi, le domaine composé ne se réduit pas aux espèces figurant dans une table métrologique.

\section{Classification interne et grammaire compositionnelle}
\label{sec:alfabeto-composicional-rev6}

Soit
\[
 \mathcal C_{81}=I_9\times I_9,
 \qquad I_9=\{1,\ldots,9\},
\]
l’atlas des cellules, et soit
\[
 \rho(r,c)=(10-r,10-c)
\]
son inversion centrale. Nous notons \(\mathcal L,\mathcal N,\mathcal D,\mathcal U\) les sous-ensembles des leptons chargés, des neutrinos, des quarks de type bas et des quarks de type haut, respectivement, et posons
\[
 \mathcal Q=\mathcal D\sqcup\mathcal U.
\]
Chaque cellule possède un niveau combinatoire \(\operatorname{gen}(x)\). Dans le secteur des quarks, le résidu de couleur est également conservé
\[
 \kappa(x)=r(x)-c(x)\pmod3.
\]
Écrivons
\[
 \mathcal Q_a=\{q\in\mathcal Q:\kappa(q)=a\},
 \qquad a\in\mathbb Z/3\mathbb Z.
\]
Par la \cref{prop:color-atlas},
\[
 |\mathcal Q_0|=|\mathcal Q_1|=|\mathcal Q_2|=12.
\]

\begin{definition}[Dual formel et identité de cellule]
\label{def:dual-identidad-composicional} Pour une section de cellule orientée \(x\), son dual formel \(x^\vee\) inverse l’orientation de charge et de couleur, et a pour support \(\rho(x)\). La notation conserve la distinction
\[
 x^\vee\ne \rho(x)
\]
lorsque \(\rho(x)\) est considérée comme cellule sans orientation. Pour chaque \(x\in\mathcal C_{81}\), nous notons \(\operatorname{id}_x:x\to x\) la flèche identité. L’ensemble des identités est
\[
 \mathfrak I
 :=\{\operatorname{id}_x:x\in\mathcal C_{81}\}.
\]
\end{definition}

L’identité \(\operatorname{id}_x\) est une flèche neutre du système compositionnel. Cette seule condition n’en fait pas une réalisation d’un boson neutre. De même, le dual formel intègre une orientation et ne peut être remplacé par une seconde cellule nue.

\begin{definition}[Expressions mésoniques et baryoniques]
\label{def:expresiones-mesonicas-barionicas} L’ensemble des expressions mésoniques neutres en couleur est
\[
 \mathfrak M
 :=
 \left\{
   M(q,q')=q\otimes{(q')}^\vee:
   q,q'\in\mathcal Q,\ \kappa(q)=\kappa(q')
 \right\}.
\]
Les paires sont ordonnées. Leur résidu total est \(\kappa(q)-\kappa(q')=0\).

L’ensemble des expressions baryoniques à positions de couleur fixées est
\[
 \mathfrak B
 :=
 \mathcal Q_0\times\mathcal Q_1\times\mathcal Q_2,
\]
et nous écrivons
\[
 B(q_0,q_1,q_2)=q_0\otimes q_1\otimes q_2.
\]
Leur neutralité résiduelle est l’identité
\[
 0+1+2=0\pmod3.
\]
\end{definition}

Ces expressions sont des objets syntaxiques typés. En particulier, aucun quotient n’a encore été pris selon les permutations, le spin, la saveur, la statistique, l’énergie de liaison ou l’équivalence dynamique.

\begin{definition}[Flèches chargées compatibles]
\label{def:flechas-cargadas-compatibles} Les flèches leptoniques orientées sont
\[
 \mathfrak A_\ell
 :=
 \left\{
 (x,y)\in
   (\mathcal L\times\mathcal N)
   \sqcup
   (\mathcal N\times\mathcal L):
 \operatorname{gen}(x)=\operatorname{gen}(y)
 \right\}.
\]
Les flèches de quarks orientées sont
\[
 \mathfrak A_q
 :=
 \left\{
 (x,y)\in
   (\mathcal U\times\mathcal D)
   \sqcup
   (\mathcal D\times\mathcal U):
 \begin{array}{l}
 \operatorname{gen}(x)=\operatorname{gen}(y),\\[-2pt]
 \kappa(x)=\kappa(y)
 \end{array}
 \right\}.
\]
Nous posons
\[
 \mathfrak A_{\rm ch}:=\mathfrak A_\ell\sqcup\mathfrak A_q.
\]
\end{definition}

L’égalité de niveau type la transition dans les deux secteurs ; l’égalité de couleur est en outre exigée dans le secteur des quarks. L’orientation distingue \(x\to y\) de \(y\to x\).

\begin{definition}[Système fini de compositions typées]
\label{def:gramatica-celular-compuesta} Le système compositionnel fini est la réunion disjointe typée
\[
 \mathfrak G_{\rm comp}
 :=
 \mathcal C_{81}
 \sqcup\mathfrak I
 \sqcup\mathfrak M
 \sqcup\mathfrak B
 \sqcup\mathfrak A_{\rm ch}.
\]
Les \(81\) cellules sont les générateurs primitifs. Les \(81\) identités sont des flèches sur ces mêmes générateurs et non un second atlas.
\end{definition}

\section{Dénombrement combinatoire des compositions de parcours indépendant des données métrologiques}
\label{sec:censo-composicional-rev6}

\begin{theorem}[Dénombrement du système fini de compositions]
\label{thm:censo-gramatica-composicional-rev6} Le système de la \cref{def:gramatica-celular-compuesta} contient
\[
 \boxed{\begin{gathered}
        81\ \text{identités},\qquad
        432\ \text{expressions mésoniques},\\
        1728\ \text{expressions baryoniques},\qquad
        372\ \text{flèches chargées}
 \end{gathered}}
\]
Les flèches chargées se décomposent exactement comme suit
\[
 372=320_{\text{leptoniques}}+52_{\text{quark}}.
\]
Le dénombrement ne dépend que de l’atlas, de son inversion, des niveaux et du résidu de couleur. Il n’utilise ni masses, ni unités, ni signatures publiées. Les quatre cardinaux comptent des expressions compositionnelles internes ou des flèches admissibles du système ; ils ne comptent ni espèces physiques distinctes, ni particules physiques, ni masses.
\end{theorem}

\begin{proof}
L’application \(x\mapsto\operatorname{id}_x\) est bijective, de sorte que
\[
 |\mathfrak I|=|\mathcal C_{81}|=81.
\]
Pour les mésons, la condition de neutralité permet de choisir indépendamment une paire ordonnée dans chaque classe de couleur. Comme chaque classe contient douze cellules,
\[
 |\mathfrak M|
 =\sum_{a\in\mathbb Z/3\mathbb Z}|\mathcal Q_a|^2
 =3\cdot12^2
 =432.
\]
Pour les baryons, les positions \(0,1,2\) sont fixées, et donc
\[
 |\mathfrak B|
 =|\mathcal Q_0||\mathcal Q_1||\mathcal Q_2|
 =12^3
 =1728.
\]

Les familles leptoniques chargées ont pour tailles
\[
 |\mathcal L_{G0}|=1,\qquad
 |\mathcal L_{G2}|=8,\qquad
 |\mathcal L_{G4}|=16,
\]
tandis que les familles neutriniques ont pour tailles
\[
 |\mathcal N_{G1}|=2,\quad
 |\mathcal N_{G2}|=4,\quad
 |\mathcal N_{G3}|=6,\quad
 |\mathcal N_{G4}|=8.
\]
Seuls \(G2\) et \(G4\) sont des niveaux communs. Comme les deux orientations sont comptées,
\[
 |\mathfrak A_\ell|
 =2(8\cdot4+16\cdot8)
 =320.
\]

Pour les flèches de quarks, dans l’ordre des couleurs \((0,1,2)\), les multiplicités de type haut et de type bas aux niveaux communs sont
\[
\begin{array}{ccc}
 &\mathcal U&\mathcal D\\
 G1&(2,1,1)&(0,1,1)\\
 G3&(4,4,4)&(2,2,2).
\end{array}
\]
Le nombre de paires d’une orientation et de même couleur est donc
\[
 (2\cdot0+1\cdot1+1\cdot1)
 +(4\cdot2+4\cdot2+4\cdot2)
 =2+24=26.
\]
En incluant l’orientation opposée, on obtient
\[
 |\mathfrak A_q|=2\cdot26=52.
\]
Enfin,
\[
 |\mathfrak A_{\rm ch}|=320+52=372.
\]
Aucune de ces règles ne contient de valeur cible ni de vecteur de masse.
\end{proof}

\begin{corollary}[Constructeur compositionnel sans cibles]
\label{cor:constructor-composicional-sin-blancos} Le constructeur qui énumère conjointement les identités, les expressions mésoniques, les expressions baryoniques et les flèches chargées se factorise par les champs typés
\[
 (\text{secteur},\text{famille},\text{niveau},\kappa,\rho)
\]
de l’atlas. Avec la notation précédente, ses quatre codomaines sont \(\mathfrak I,\mathfrak M,\mathfrak B,\mathfrak A_{\rm ch}\). En particulier, le dénombrement peut être fixé avant l’ouverture de toute table de signatures.
\end{corollary}

\begin{proof}
Les quatre définitions et le calcul du \cref{thm:censo-gramatica-composicional-rev6} n’emploient que ces champs. La table de signatures n’appartient au domaine d’aucune des règles.
\end{proof}

L’indépendance à l’égard des cibles ne signifie pas que les expressions ont déjà acquis un nom physique. Elle signifie quelque chose d’antérieur et de plus précis : il existe un domaine compositionnel construit sans utiliser comme critère la quantité qui sera évaluée ensuite.

\section{Transport de mémoire du parcours composé à la signature entière}
\label{sec:memoria-composicion-especies-rev6}

Soit \(\gamma\in\Gamma_{108}^{\rm enr}\) une présentation enrichie d’une extension. Son enregistrement lu contient
\[
 \mathcal E(\gamma)
 =
 (n_A,e_0,\ldots,e_{11},T_\zeta,C_\zeta),
\]
outre les données de phase, de feuillet, d’orientation et de bord qui typent une composition. Les douze événements ne peuvent être remplacés par leur total avant la composition. Pour \(0\le i<6\), nous définissons
\[
 p_i=e_i+e_{i+6},
 \qquad
 a_i=e_i-e_{i+6},
\]
\[
 s_C=\sum_{i=0}^{5}p_i,
 \qquad
 M_i=p_i-p_5\quad(0\le i<5).
\]
Avec \(M_5=(M_0,\ldots,M_4)\) et \(A_6=(a_0,\ldots,a_5)\), la reconstruction est
\[
 p_5=\frac{s_C-\sum_{i=0}^{4}M_i}{6},
 \qquad
 p_i=p_5+M_i,
\]
\[
 e_i=\frac{p_i+a_i}{2},
 \qquad
 e_{i+6}=\frac{p_i-a_i}{2}.
\]
D’après le \cref{thm:memoria-1-5-6}, on obtient sur l’image entière la décomposition réversible
\[
 \boxed{(s_C,M_5,A_6),\qquad 1+5+6=12.}
\]
Le premier bloc conserve le total visible ; le deuxième, la forme relative dans le feuillet symétrique ; le troisième, la mémoire entre les deux demi-couronnes.

La coordonnée torsionnelle forte requiert en outre la cochaîne au niveau du pas :
\[
 k_{\Delta_4}=T_\zeta-2C_\zeta.
\]
Elle n’est pas une fonction du seul motif d’événements nuls. Cette différence sera essentielle pour interpréter le losange de non-descente.

\begin{definition}[Donnée de composition orientée]
\label{def:dato-empalme-orientado-rev6} Soient \(\Gamma,\Lambda\) deux enregistrements enrichis. Une donnée de composition est une paire \((g,B)\), où \(g\) aligne la phase, le feuillet, l’orientation et la position à l’interface, et \(B\) est la sous-trace commune, avec sa contribution à l’action, ses douze événements, ses prédécesseurs torsionnels et sa couronne. Sur le domaine où ces données sont compatibles, on définit
\[
 \boxed{
 {(n,e)}_\Gamma\star_{(g,B)}{(n,e)}_\Lambda
 =
 \bigl(
 n_\Gamma+n_\Lambda-n_B,\,
 e_\Gamma+g e_\Lambda-e_B
 \bigr).}
\]
Le mot torsionnel et la couronne se composent au même niveau, avant le calcul de \(T_\zeta,C_\zeta\) et de \(k_{\Delta_4}\).
\end{definition}

La formule est une inclusion-exclusion d’histoires sur une sous-trace identifiée. Elle est démontrée avec une structure initiale explicite : les enregistrements, le transport \(g\), le bord \(B\), l’orientation et le mot des prédécesseurs font partie des données d’entrée. Le système compositionnel ne choisit de lui-même aucune de ces données.

\begin{definition}[Extracteur postérieur à la composition]
\label{def:extractor-posterior-empalme-rev6} Pour un enregistrement enrichi complet, on définit
\[
 L_{\rm tick}(\gamma)
 =
 \bigl(
 n_A,s_C,k_{\Delta_4},
 \nu_{120}^{\rm ev},\nu_{270}
 \bigr)\in\ZZ^5,
\]
où, si \(\tau_m=\operatorname{sgn}(e_m)\) et \(I_a=\{a,a+4,a+8\}\pmod {12}\),
\[
 \nu_{120}^{\rm ev}
 =
 \sum_{a=0}^{3}
 \operatorname{sgn}\left(\sum_{m\in I_a}e_m\right),
 \qquad
 \nu_{270}
 =
 \sum_{m=0}^{11}\tau_m\tau_{m+3}.
\]
La signature d’une composition est définie par
\[
 \operatorname{Sig}(\Gamma,\Lambda;g,B)
 :=
 L_{\rm tick}
 \bigl(\Gamma\star_{(g,B)}\Lambda\bigr).
\]
\end{definition}

Les signes et les autocorrélations s’appliquent à l’enregistrement déjà composé. La définition n’équivaut donc pas à additionner les signatures des constituants.

\section{Diagramme commutatif de composition et de préservation de la fibre}
\label{sec:rombo-no-descenso-composicional-rev6}

Le catalogue historique B1 emploie une projection d’événements qu’il convient de typer séparément. Soit
\[
 \Sigma_{12}=(+,+,+,-,-,-,+,+,+,-,-,-),
 \qquad \tau_m=\operatorname{sgn}(e_m),
\]
et définissons le proxy d’événements nuls
\[
 K_0(e)
 =
 \sum_{\tau_m=0}\Sigma_{12}(m).
\]
La projection exacte est
\[
 F_0(n_A,e)
 =
 \left(
 n_A,\,
 \sum_m e_m,\,
 K_0(e),\,
 \sum_{a=0}^{3}
   \operatorname{sgn}\sum_{m\in I_a}e_m,\,
 \sum_m\tau_m\tau_{m+3}
 \right)\in\ZZ^5.
\]
Pour deux enregistrements alignés, nous écrivons
\[
 (n,e)\boxplus(n',e')=(n+n',e+e').
\]
Cette opération est le cas sans sous-trace commune et avec alignement identité de la composition d’événements.

\begin{proposition}[Obstruction à la descente de la projection B1]
\label{prop:no-descenso-proyeccion-b1-rev6} Il n’existe aucune opération binaire
\[
 \mu:\ZZ^5\times\ZZ^5\longrightarrow\ZZ^5
\]
tel que
\[
 F_0(\Gamma\boxplus\Lambda)
 =
 \mu\bigl(F_0(\Gamma),F_0(\Lambda)\bigr)
\]
pour tous les enregistrements B1.
\end{proposition}

\begin{proof}
Considérons les parcours, écrits dans la carte \((r,c,h_0,v_0)\),
\[
 a=(1,3,-1,+1),
 \qquad
 b=(6,7,-1,+1),
 \qquad
 c=(1,2,-1,-1).
\]
Les deux premiers ont des enregistrements d’événements distincts :
\[
\begin{aligned}
 e(a)&=(5,4,4,-4,-4,-6,\;5,4,4,-4,-4,-6),\\
 e(b)&=(4,5,4,-4,-6,-4,\;4,5,4,-4,-6,-4),
\end{aligned}
\]
mais produisent le même quintuplet :
\[
 F_0(a)=F_0(b)=(8,-2,0,0,-12).
\]
Le parcours de test possède
\[
 e(c)=(1,4,1,-2,-4,-3,\;1,4,1,-2,-4,-3)
\]
et
\[
 F_0(c)=(28,-6,0,-4,-12).
\]
En additionnant d’abord les événements puis en réévaluant, on obtient
\[
 F_0(a\boxplus c)=(36,-8,0,0,-12),
\]
\[
 F_0(b\boxplus c)=(36,-8,0,-2,-12).
\]
Les deux paires d’entrées de l’hypothétique \(\mu\) sont égales, mais les sorties diffèrent en \(\nu_{120}^{\rm ev}\). Cela contredit l’existence de \(\mu\).
\end{proof}

Le dénombrement exhaustif du même domaine contient \(324\) parcours, \(53\) quintuplets proxy et \(198\) états d’événements distincts. Quarante-deux fibres de la projection contiennent plus d’un état enrichi ; quarante d’entre elles sont séparées par un parcours de test. Le losange précédent est le témoin canonique le plus bref de cette perte d’information.

\begin{remark}[Séparation typée entre le représentant \(K_0\) et la torsion forte]
\label{rem:proxy-no-torsion-fuerte-rev6} La troisième coordonnée de \(F_0\) est \(K_0(e)\), qui, dans B1, prend ses valeurs dans \(\{-4,-2,0,2,4\}\). Elle n’est pas identifiée à \(k_{\Delta_4}=T_\zeta-2C_\zeta\). Il existe en effet des parcours de même \(n_A\) et de même vecteur de douze événements auxquels le lecteur torsionnel attribue des valeurs distinctes parce qu’il conserve les prédécesseurs de chaque itération et la couronne.

La \cref{prop:no-descenso-proyeccion-b1-rev6} démontre précisément que la projection B1 à cinq coordonnées ne transporte pas la superposition des enregistrements. Elle n’affirme aucune impossibilité universelle pour toute structure future sur \(\ZZ^5\) et n’autorise pas à remplacer \(K_0\) par \(k_{\Delta_4}\). Dans le lecteur fort, la sortie est définie par la \cref{def:extractor-posterior-empalme-rev6} : on compose d’abord le mot enrichi, puis on en extrait à nouveau les cinq coordonnées.
\end{remark}

\section{Théorème de composition des enregistrements indépendant des données métrologiques}
\label{sec:teorema-composicional-sin-blancos-rev6}

\begin{theorem}[Clôture compositionnelle sans cibles]
\label{thm:cierre-composicional-sin-blancos-rev6} Une fois fixés l’atlas \(\mathcal C_{81}\), la dualité formelle, le résidu de couleur et l’enregistrement enrichi du \cref{chap:extension-ruta-caracter}, les affirmations suivantes sont vérifiées.
\begin{enumerate}
 \item Il existe le système fini de compositions typées \(\mathfrak G_{\rm comp}\) avec \(81\) identités, \(432\) expressions mésoniques, \(1728\) expressions baryoniques et \(372\) flèches chargées, décomposées comme \(320+52\).

 \item Le système est construit sans masses cibles. Sa sortie est une expression, une orbite ou une multisection typée, et non une masse ni une étiquette physique individualisée.

 \item Étant donnés une expression et des enregistrements constituants munis d’un arbre explicite de compositions compatibles \((g,B)\), l’opération produit un enregistrement enrichi. Sa signature est obtenue par l’application
 \[
 \boxed{
 \bigl(
 \text{expression},
 \text{registres},
 \text{compositions orientées}
 \bigr)
 \xrightarrow{\ \mathfrak C_{12}\ }
 \text{registre composé}
 \xrightarrow{\ L_{\rm tick}\ }
 \ZZ^5.}
 \]
 Ici, \(\mathfrak C_{12}\) désigne l’itération de \(\star_{(g,B)}\) avec le parenthésage et les bords inclus dans les données.

 \item La projection B1 vers cinq entiers ne constitue pas à elle seule un espace compositionnel : le diagramme avec \(\boxplus\) ne descend pas à \(\ZZ^5\). Par conséquent, remplacer chaque enregistrement par son quintuplet avant la composition fait perdre une information nécessaire à la réévaluation d’au moins \(\nu_{120}^{\rm ev}\).
\end{enumerate}
\end{theorem}

\begin{proof}
Le premier point est le \cref{thm:censo-gramatica-composicional-rev6}. Le deuxième est le \cref{cor:constructor-composicional-sin-blancos}. Pour le troisième, chaque arête de l’arbre de composition est évaluée au moyen de la \cref{def:dato-empalme-orientado-rev6} ; comme les bords et les transports font partie du domaine, l’itération produit un enregistrement sur lequel \(L_{\rm tick}\) est défini. Aucune signature intermédiaire ne remplace l’enregistrement. Le quatrième point est la \cref{prop:no-descenso-proyeccion-b1-rev6}.
\end{proof}

La formule correcte de la signature composée est donc
\[
 \operatorname{Sig}_{\rm comp}
 =
 L_{\rm tick}\circ\mathfrak C_{12}.
\]
Il ne s’agit pas d’une somme de lignes de masse. Ce n’est qu’ensuite que le caractère formel ou sa spécialisation globale peuvent être appliqués :
\[
 \text{registre composé}
 \xrightarrow{L_{\rm tick}}\ZZ^5
 \xrightarrow{\chi_{\rm form}}\operatorname{Fun}(\mathcal P,\RR_{>0})
 \longrightarrow
 \text{carte dimensionnelle}.
\]

\section{Reconstruction généalogique du mélange, de la génération et de la persistance de saveur}
\label{sec:significado-genealogico-composicion-rev6}

Le théorème précédent complète une strate qu’une liste de masses ne peut montrer. Les cellules primitives ne se limitent pas à recevoir des noms ; elles engendrent des identités, des paires duales, des triplets de couleur et des flèches orientées selon des règles internes. Les objets composés apparaissent ainsi dans le domaine avant toute comparaison métrologique.

La généalogie ne se réduit pas non plus à l’arbre syntaxique. Deux expressions comportant le même nombre de constituants peuvent se composer différemment parce que la phase, le feuillet, l’orientation et la sous-trace commune font partie de la donnée. L’enregistrement conserve ces différences jusqu’à ce que les signes, les corrélations et la torsion aient agi. La masse, lorsqu’une réalisation dimensionnelle existe, sera une lecture ultérieure de cette mémoire composée.

Cet ordre explique pourquoi \(\ZZ^5\) n’est pas le domaine de composition, bien qu’il soit celui du caractère. Le groupe \((\ZZ^5,+)\) existe et le caractère y est multiplicatif ; ce qui n’existe pas pour la projection B1, c’est une loi qui récupère à partir des quintuplets la superposition ayant déjà oublié la disposition des douze événements. La linéarité du caractère ne restitue pas l’information supprimée par l’extracteur.

\section{Du système compositionnel interne à la réalisation physique}
\label{sec:delimitacion-composicion-especies-rev6}

Le \cref{thm:cierre-composicional-sin-blancos-rev6} construit le domaine antérieur aux noms physiques : expressions ordonnées, duaux, triplets de couleur et flèches compatibles. Ses cardinaux comptent des possibilités internes de composition. Une espèce physique, en revanche, est une représentation persistante dotée d’une saveur, d’une couleur, d’un spin, d’une statistique et d’un canal déterminés. Confondre ces deux niveaux transformerait un dénombrement d’expressions en un faux catalogue de particules.

\begin{definition}[Donnée de réalisation physique]
\label{def:dato-realizacion-fisica-rev6} Une donnée de réalisation pour une espèce \(X\) est le tuple
\[
 \mathfrak D_X=
 \bigl(
 \mathfrak t_X,\mathfrak s_X,
 (\mathcal L_{X_i})_{i=1}^{r},
 \mathcal T_X,(g_X,B_X)
 \bigr),
\]
où :
\begin{enumerate}
 \item \(\mathfrak t_X\) fixe le type physique --- élémentaire, nul, mésonique, baryonique, multiquark ou topologique --- et sa représentation ;
 \item \(\mathfrak s_X\) est la multisection ou l’expression interne réalisée ;
 \item \(\mathcal L_{X_i}\) sont les enregistrements enrichis des constituants ;
 \item \(\mathcal T_X\) fixe la saveur, la couleur, le spin, la statistique et la symétrisation ;
 \item \((g_X,B_X)\) fixe l’arbre de composition, ses transports et ses sous-traces communes.
\end{enumerate}
Le tuple est déclaré sans consulter la masse ou la largeur qui seront évaluées ensuite.
\end{definition}

\begin{proposition}[La réalisation modifie l’enregistrement, non la loi]
\label{prop:realizacion-compuesta-ley-unica-rev6} Toute donnée \(\mathfrak D_X\) détermine un enregistrement composé
\[
 \mathcal L_X
 =\mathfrak C_{12}
   \bigl((\mathcal L_{X_i})_{i=1}^{r};
         \mathcal T_X,g_X,B_X\bigr),
\]
une signature \(v_X=L_{\rm tick}(\mathcal L_X)\) et le caractère central universel
\[
 \chi_0(v_X-v_e).
\]
Si la généalogie apporte en outre un raffinement spectral typé \(\eta_X\in\mathcal E_{\mathfrak S}\), l’évaluation complète est
\[
 m_X
 =m_e^{\rm HMT}\,
   \chi_{\rm full}(v_X-v_e,\eta_X).
\]
Le caractère central et le raffinement appartiennent à des niveaux distincts ; la complétude de la tour ne remplace pas l’application prospective produisant \(\eta_X\). Si \(X\) est instable, sa largeur relève du lecteur temporel \(\Gamma_X=\hbar_{\rm int}/\tau_X^{\rm life}\), et non d’une sixième coordonnée de la signature de masse.
\end{proposition}

\begin{proof}
La donnée de réalisation contient exactement les entrées de \(\mathfrak C_{12}\). Le \cref{thm:cierre-composicional-sin-blancos-rev6} produit l’enregistrement enrichi et la \cref{def:extractor-posterior-empalme-rev6} produit sa signature. Le caractère central et son relèvement spectral sont ceux de \cref{eq:caracter-masa-global-rev7,eq:ley-masa-todos-registros-rev6}. La largeur utilise une droite temporelle et possède donc un autre domaine.
\end{proof}

La séparation des niveaux est ainsi fixée :
\[
\boxed{
\begin{gathered}
\text{système compositionnel interne}
\longrightarrow\text{donnée de réalisation}
\longrightarrow\text{registre composé}\\
\longrightarrow\text{signature}
\longrightarrow\text{caractère central}
\longrightarrow\text{raffinement spectral}
\longrightarrow\text{comparaison physique}.
\end{gathered}}
\]
Les \(432\) expressions mésoniques, les \(1728\) expressions baryoniques et les \(372\) flèches chargées n’occupent que la première case et ne constituent pas un inventaire de particules. Le Modèle standard conserve son ensemble typé d’espèces, et chaque espèce occupe la chaîne au moyen de sa propre donnée de réalisation. La règle causale est unique dans tous les secteurs :
\[
 \boxed{\text{composer les extensions et les registres ;\quad lire ensuite la signature}.}
\]
